\documentclass[12pt,a4paper]{article}
\usepackage[UTF8]{ctex}
\usepackage{amsmath,amssymb,amsthm}
\usepackage{geometry}

\geometry{margin=2.5cm}

\title{刚体速度合成定理}
\author{第8章补充说明}
\date{}

\begin{document}

\maketitle

\section{问题提出}

\textbf{问题}：如何理解角速度和线速度的合成公式？它们的官方名称是什么？

\textbf{公式}：

\subsection{角速度合成公式}

如果坐标系$\{i\}$相对于坐标系$\{i-1\}$有角速度$\omega_{\text{rel}}$，那么：
\begin{equation}
\omega_i = R_i^T \omega_{i-1} + \omega_{\text{rel}} \tag{3}
\end{equation}

\subsection{线速度合成公式}

如果坐标系$\{i\}$的原点相对于坐标系$\{i-1\}$有速度，那么：
\begin{equation}
v_i = R_i^T (v_{i-1} + \omega_{i-1} \times p_i) + v_{\text{rel}} \tag{4}
\end{equation}

其中$v_{\text{rel}}$是坐标系$\{i\}$相对于坐标系$\{i-1\}$的相对速度。

\section{一般形式}

\subsection{角速度合成定理的一般形式}

\textbf{最一般形式}：

设坐标系$\{a\}$和$\{b\}$，其中$\{b\}$相对于$\{a\}$有相对运动。那么：
\begin{equation}
\omega_b = R_{ba}^T \omega_a + \omega_{\text{rel}} \tag{3-一般}
\end{equation}

\textbf{其中}：
\begin{itemize}
\item $\omega_b \in \mathbb{R}^3$：坐标系$\{b\}$的绝对角速度（在$\{b\}$中表示）
\item $\omega_a \in \mathbb{R}^3$：坐标系$\{a\}$的绝对角速度（在$\{a\}$中表示）
\item $R_{ba} \in SO(3)$：从坐标系$\{a\}$到坐标系$\{b\}$的旋转矩阵
\item $\omega_{\text{rel}} \in \mathbb{R}^3$：坐标系$\{b\}$相对于$\{a\}$的相对角速度（在$\{b\}$中表示）
\end{itemize}

\textbf{特殊情况}：
\begin{enumerate}
\item \textbf{纯旋转关节}：$\omega_{\text{rel}} = \dot{\theta}_i z_i$，其中$z_i$是旋转轴（在$\{i\}$中表示）
\item \textbf{纯移动关节}：$\omega_{\text{rel}} = 0$
\item \textbf{螺旋关节}：$\omega_{\text{rel}} = \dot{\theta}_i z_i$（与线速度耦合）
\item \textbf{一般关节}：$\omega_{\text{rel}}$可以是任意向量（在$\{i\}$中表示）
\end{enumerate}

\subsection{线速度合成定理的一般形式}

\textbf{最一般形式}：

设坐标系$\{a\}$和$\{b\}$，其中$\{b\}$相对于$\{a\}$有相对运动。那么：
\begin{equation}
v_b = R_{ba}^T (v_a + \omega_a \times p_{ab}) + v_{\text{rel}} \tag{4-一般}
\end{equation}

\textbf{其中}：
\begin{itemize}
\item $v_b \in \mathbb{R}^3$：坐标系$\{b\}$原点的绝对线速度（在$\{b\}$中表示）
\item $v_a \in \mathbb{R}^3$：坐标系$\{a\}$原点的绝对线速度（在$\{a\}$中表示）
\item $\omega_a \in \mathbb{R}^3$：坐标系$\{a\}$的绝对角速度（在$\{a\}$中表示）
\item $p_{ab} \in \mathbb{R}^3$：从坐标系$\{a\}$原点到坐标系$\{b\}$原点的向量（在$\{a\}$中表示）
\item $R_{ba} \in SO(3)$：从坐标系$\{a\}$到坐标系$\{b\}$的旋转矩阵
\item $v_{\text{rel}} \in \mathbb{R}^3$：坐标系$\{b\}$相对于$\{a\}$的相对线速度（在$\{b\}$中表示）
\end{itemize}

\textbf{特殊情况}：
\begin{enumerate}
\item \textbf{纯旋转关节}：$v_{\text{rel}} = 0$
\item \textbf{纯移动关节}：$v_{\text{rel}} = \dot{d}_i z_i$，其中$z_i$是移动方向（在$\{i\}$中表示）
\item \textbf{螺旋关节}：$v_{\text{rel}} = h \dot{\theta}_i z_i$，其中$h$是螺旋节距
\item \textbf{一般关节}：$v_{\text{rel}}$可以是任意向量（在$\{i\}$中表示）
\end{enumerate}

\subsection{用Twist（螺旋运动）表示的统一形式}

\textbf{Twist定义}：

将角速度和线速度组合成6维twist：
\begin{equation}
V_b = \begin{bmatrix} \omega_b \\ v_b \end{bmatrix} \in \mathbb{R}^6
\end{equation}

\textbf{统一形式}：
\begin{equation}
V_b = \text{Ad}_{T_{ba}^{-1}}(V_a) + A_{\text{rel}} \dot{\theta} \tag{统一形式}
\end{equation}

\textbf{其中}：
\begin{itemize}
\item $V_b = [\omega_b^T, v_b^T]^T$：坐标系$\{b\}$的twist（在$\{b\}$中表示）
\item $V_a = [\omega_a^T, v_a^T]^T$：坐标系$\{a\}$的twist（在$\{a\}$中表示）
\item $T_{ba} = \begin{bmatrix} R_{ba} & p_{ab} \\ 0 & 1 \end{bmatrix}$：从$\{a\}$到$\{b\}$的齐次变换矩阵
\item $\text{Ad}_{T_{ba}^{-1}}$：伴随变换（Adjoint transformation）
\item $A_{\text{rel}} = [\omega_{\text{rel}}^T, v_{\text{rel}}^T]^T$：相对运动的螺旋轴（screw axis，在$\{b\}$中表示）
\item $\dot{\theta}$：相对运动的速度参数
\end{itemize}

\textbf{伴随变换的展开形式}：
\begin{equation}
\text{Ad}_{T_{ba}^{-1}} = \begin{bmatrix}
R_{ba}^T & 0 \\
-[R_{ba}^T p_{ab}]_{\times} R_{ba}^T & R_{ba}^T
\end{bmatrix}
\end{equation}

其中$[p]_{\times}$是$p$的反对称矩阵。

\textbf{展开后的形式}：
\begin{align}
\begin{bmatrix} \omega_b \\ v_b \end{bmatrix} &= \begin{bmatrix}
R_{ba}^T & 0 \\
-[R_{ba}^T p_{ab}]_{\times} R_{ba}^T & R_{ba}^T
\end{bmatrix} \begin{bmatrix} \omega_a \\ v_a \end{bmatrix} + A_{\text{rel}} \dot{\theta} \\
&= \begin{bmatrix}
R_{ba}^T \omega_a \\
-[R_{ba}^T p_{ab}]_{\times} R_{ba}^T \omega_a + R_{ba}^T v_a
\end{bmatrix} + A_{\text{rel}} \dot{\theta}
\end{align}

\textbf{验证}：

展开后得到：
\begin{align}
\omega_b &= R_{ba}^T \omega_a + \omega_{\text{rel}} \\
v_b &= -[R_{ba}^T p_{ab}]_{\times} R_{ba}^T \omega_a + R_{ba}^T v_a + v_{\text{rel}}
\end{align}

注意到：
\begin{equation}
-[R_{ba}^T p_{ab}]_{\times} R_{ba}^T \omega_a = R_{ba}^T (p_{ab} \times \omega_a) = R_{ba}^T (\omega_a \times (-p_{ab})) = R_{ba}^T (\omega_a \times p_{ab})
\end{equation}

因此：
\begin{equation}
v_b = R_{ba}^T (v_a + \omega_a \times p_{ab}) + v_{\text{rel}}
\end{equation}

这与公式(4-一般)一致。

\subsection{矩阵形式的完整表示}

\textbf{完整形式}：
\begin{equation}
\begin{bmatrix} \omega_b \\ v_b \end{bmatrix} = \begin{bmatrix}
R_{ba}^T & 0 \\
R_{ba}^T [p_{ab}]_{\times} & R_{ba}^T
\end{bmatrix} \begin{bmatrix} \omega_a \\ v_a \end{bmatrix} + \begin{bmatrix} \omega_{\text{rel}} \\ v_{\text{rel}} \end{bmatrix} \tag{矩阵形式}
\end{equation}

\textbf{其中}：
\begin{itemize}
\item $[p_{ab}]_{\times}$是$p_{ab}$的反对称矩阵：
\begin{equation}
[p_{ab}]_{\times} = \begin{bmatrix}
0 & -p_{ab,z} & p_{ab,y} \\
p_{ab,z} & 0 & -p_{ab,x} \\
-p_{ab,y} & p_{ab,x} & 0
\end{bmatrix}
\end{equation}
\item $\omega_{\text{rel}}$和$v_{\text{rel}}$都在坐标系$\{b\}$中表示
\end{itemize}

\textbf{物理意义}：
\begin{itemize}
\item 第一行：角速度合成（公式3-一般）
\item 第二行：线速度合成（公式4-一般），其中$R_{ba}^T [p_{ab}]_{\times} \omega_a = R_{ba}^T (\omega_a \times p_{ab})$
\end{itemize}

\textbf{符号说明}：
\begin{itemize}
\item $\{a\}$：参考坐标系（或前一个坐标系）
\item $\{b\}$：当前坐标系（或后一个坐标系）
\item $R_{ba}$：从$\{a\}$到$\{b\}$的旋转矩阵
\item $p_{ab}$：从$\{a\}$原点到$\{b\}$原点的向量（在$\{a\}$中表示）
\end{itemize}

\section{官方名称}

\subsection{中文名称}

\begin{itemize}
\item \textbf{角速度合成定理}（Angular Velocity Composition Theorem）
\item \textbf{线速度合成定理}（Linear Velocity Composition Theorem）
\item 或统称为\textbf{刚体速度合成定理}（Rigid Body Velocity Composition Theorem）
\end{itemize}

\subsection{英文名称}

\begin{itemize}
\item \textbf{Angular Velocity Composition Theorem}
\item \textbf{Linear Velocity Composition Theorem}
\item 或统称为\textbf{Rigid Body Velocity Composition Theorem}
\item 也称为\textbf{Velocity Addition Theorem}（速度叠加定理）
\end{itemize}

\subsection{在经典力学中的名称}

在经典力学中，这些公式通常被称为：
\begin{itemize}
\item \textbf{绝对速度 = 牵连速度 + 相对速度}
\item \textbf{绝对角速度 = 牵连角速度 + 相对角速度}
\end{itemize}

其中：
\begin{itemize}
\item \textbf{绝对速度/角速度}：物体相对于固定参考系的速度/角速度
\item \textbf{牵连速度/角速度}：由于参考坐标系运动产生的速度/角速度
\item \textbf{相对速度/角速度}：物体相对于参考坐标系的速度/角速度
\end{itemize}

\section{物理意义}

\subsection{角速度合成公式（公式3）}

\textbf{物理意义}：

坐标系$\{i\}$的绝对角速度$\omega_i$由两部分组成：
\begin{enumerate}
\item \textbf{牵连角速度}：$R_i^T \omega_{i-1}$
\begin{itemize}
\item 这是坐标系$\{i-1\}$的角速度$\omega_{i-1}$传递到坐标系$\{i\}$中的部分
\item 需要坐标变换（左乘$R_i^T$）从$\{i-1\}$转换到$\{i\}$
\end{itemize}

\item \textbf{相对角速度}：$\omega_{\text{rel}}$
\begin{itemize}
\item 这是坐标系$\{i\}$相对于坐标系$\{i-1\}$的角速度
\item \textbf{表示约定}：在机器人学中，通常$\omega_{\text{rel}}$在坐标系$\{i\}$中表示
\item \textbf{原因}：关节轴$z_i$通常在坐标系$\{i\}$中定义，因此$\omega_{\text{rel}} = \dot{\theta}_i z_i$也在$\{i\}$中表示
\end{itemize}
\end{enumerate}

\textbf{数学表示}：
\begin{equation}
\omega_i = \omega_{\text{牵连}} + \omega_{\text{相对}} = R_i^T \omega_{i-1} + \omega_{\text{rel}}
\end{equation}

\subsection{线速度合成公式（公式4）}

\textbf{物理意义}：

坐标系$\{i\}$原点的绝对线速度$v_i$由三部分组成：
\begin{enumerate}
\item \textbf{牵连速度的平移部分}：$R_i^T v_{i-1}$
\begin{itemize}
\item 这是坐标系$\{i-1\}$原点的速度$v_{i-1}$传递到坐标系$\{i\}$中的部分
\item 需要坐标变换（左乘$R_i^T$）从$\{i-1\}$转换到$\{i\}$
\end{itemize}

\item \textbf{牵连速度的旋转部分}：$R_i^T (\omega_{i-1} \times p_i)$
\begin{itemize}
\item 这是由于坐标系$\{i-1\}$的旋转产生的速度
\item $p_i$是从坐标系$\{i-1\}$原点到坐标系$\{i\}$原点的向量
\item $\omega_{i-1} \times p_i$是点$O_i$由于$\{i-1\}$旋转而产生的速度（在$\{i-1\}$中表示）
\item 需要坐标变换（左乘$R_i^T$）转换到$\{i\}$
\end{itemize}

\item \textbf{相对速度}：$v_{\text{rel}}$
\begin{itemize}
\item 这是坐标系$\{i\}$相对于坐标系$\{i-1\}$的相对速度
\item \textbf{表示约定}：在机器人学中，通常$v_{\text{rel}}$在坐标系$\{i\}$中表示
\item \textbf{原因}：关节移动方向$z_i$通常在坐标系$\{i\}$中定义，因此$v_{\text{rel}} = \dot{d}_i z_i$也在$\{i\}$中表示
\end{itemize}
\end{enumerate}

\textbf{数学表示}：
\begin{equation}
v_i = v_{\text{牵连}} + v_{\text{相对}} = R_i^T (v_{i-1} + \omega_{i-1} \times p_i) + v_{\text{rel}}
\end{equation}

\section{详细推导}

\subsection{角速度合成公式的推导}

\textbf{基本原理}：

角速度是矢量，可以按照矢量加法进行合成。

\textbf{步骤1：牵连角速度}

坐标系$\{i-1\}$的角速度$\omega_{i-1}$是在坐标系$\{i-1\}$中表示的。要得到它在坐标系$\{i\}$中的表示，需要进行坐标变换：
\begin{equation}
\omega_{\text{牵连}} = R_i^T \omega_{i-1}
\end{equation}

\textbf{步骤2：相对角速度}

坐标系$\{i\}$相对于坐标系$\{i-1\}$的角速度$\omega_{\text{rel}}$已经在坐标系$\{i\}$中表示（或可以转换到$\{i\}$中）。

\textbf{步骤3：合成}

根据矢量加法，坐标系$\{i\}$的绝对角速度为：
\begin{equation}
\omega_i = \omega_{\text{牵连}} + \omega_{\text{相对}} = R_i^T \omega_{i-1} + \omega_{\text{rel}} \tag{3}
\end{equation}

\subsection{线速度合成公式的推导}

\textbf{基本原理}：

刚体上任意点的速度 = 基点速度 + 由于旋转产生的速度 + 相对速度

\textbf{步骤1：基点速度}

坐标系$\{i-1\}$原点的速度$v_{i-1}$是在坐标系$\{i-1\}$中表示的。要得到它在坐标系$\{i\}$中的表示，需要进行坐标变换：
\begin{equation}
v_{\text{基点}} = R_i^T v_{i-1}
\end{equation}

\textbf{步骤2：旋转产生的速度}

考虑坐标系$\{i\}$的原点$O_i$，它在坐标系$\{i-1\}$中的位置为$p_i$。

由于坐标系$\{i-1\}$有角速度$\omega_{i-1}$，点$O_i$由于旋转而产生的速度为：
\begin{equation}
v_{\text{旋转}} = \omega_{i-1} \times p_i
\end{equation}

（在坐标系$\{i-1\}$中表示）

要转换到坐标系$\{i\}$中，需要左乘$R_i^T$：
\begin{equation}
v_{\text{旋转},\{i\}} = R_i^T (\omega_{i-1} \times p_i)
\end{equation}

\textbf{步骤3：相对速度}

坐标系$\{i\}$相对于坐标系$\{i-1\}$的相对速度$v_{\text{rel}}$。

\textbf{表示约定}：在机器人学中，$v_{\text{rel}}$通常在坐标系$\{i\}$中表示，因为：
\begin{itemize}
\item 关节移动方向$z_i$在坐标系$\{i\}$中定义
\item 因此$v_{\text{rel}} = \dot{d}_i z_i$也在$\{i\}$中表示
\end{itemize}

\textbf{步骤4：合成}

根据速度合成定理，坐标系$\{i\}$原点的绝对速度为：
\begin{equation}
v_i = v_{\text{基点}} + v_{\text{旋转}} + v_{\text{相对}} = R_i^T v_{i-1} + R_i^T (\omega_{i-1} \times p_i) + v_{\text{rel}}
\end{equation}

整理得到：
\begin{equation}
v_i = R_i^T (v_{i-1} + \omega_{i-1} \times p_i) + v_{\text{rel}} \tag{4}
\end{equation}

\section{在机器人学中的应用}

\subsection{旋转关节}

对于旋转关节$i$：
\begin{itemize}
\item $\omega_{\text{rel}} = \dot{\theta}_i z_i$（关节绕轴$z_i$旋转）
\item $v_{\text{rel}} = 0$（关节只旋转，不移动）
\end{itemize}

因此：
\begin{align}
\omega_i &= R_i^T \omega_{i-1} + \dot{\theta}_i z_i \\
v_i &= R_i^T (v_{i-1} + \omega_{i-1} \times p_i)
\end{align}

\subsection{移动关节}

对于移动关节$i$：
\begin{itemize}
\item $\omega_{\text{rel}} = 0$（关节只移动，不旋转）
\item $v_{\text{rel}} = \dot{d}_i z_i$（关节沿轴$z_i$移动）
\end{itemize}

因此：
\begin{align}
\omega_i &= R_i^T \omega_{i-1} \\
v_i &= R_i^T (v_{i-1} + \omega_{i-1} \times p_i) + \dot{d}_i z_i
\end{align}

\subsection{螺旋关节（既有$\omega_{\text{rel}}$和$v_{\text{rel}}$）}

对于螺旋关节（screw joint）或一般关节，坐标系$\{i\}$相对于坐标系$\{i-1\}$既有旋转又有移动：

\textbf{一般情况}：
\begin{itemize}
\item $\omega_{\text{rel}} \neq 0$（有相对旋转）
\item $v_{\text{rel}} \neq 0$（有相对移动）
\end{itemize}

\textbf{速度合成公式}：
\begin{align}
\omega_i &= R_i^T \omega_{i-1} + \omega_{\text{rel}} \tag{3} \\
v_i &= R_i^T (v_{i-1} + \omega_{i-1} \times p_i) + v_{\text{rel}} \tag{4}
\end{align}

\textbf{物理意义}：
\begin{itemize}
\item 角速度$\omega_i$由两部分组成：
\begin{enumerate}
\item 牵连角速度：$R_i^T \omega_{i-1}$（从$\{i-1\}$传递过来的角速度）
\item 相对角速度：$\omega_{\text{rel}}$（$\{i\}$相对于$\{i-1\}$的旋转）
\end{enumerate}

\item 线速度$v_i$由三部分组成：
\begin{enumerate}
\item 牵连速度的平移部分：$R_i^T v_{i-1}$（从$\{i-1\}$传递过来的平移速度）
\item 牵连速度的旋转部分：$R_i^T (\omega_{i-1} \times p_i)$（由于$\{i-1\}$旋转产生的速度）
\item 相对速度：$v_{\text{rel}}$（$\{i\}$相对于$\{i-1\}$的移动）
\end{enumerate}
\end{itemize}

\textbf{特殊情况：螺旋运动}

如果相对旋转和相对移动沿着同一轴（螺旋轴），那么：
\begin{itemize}
\item $\omega_{\text{rel}} = \dot{\theta}_i z_i$（绕轴$z_i$旋转）
\item $v_{\text{rel}} = h \dot{\theta}_i z_i$（沿轴$z_i$移动，$h$是螺旋节距）
\end{itemize}

因此：
\begin{align}
\omega_i &= R_i^T \omega_{i-1} + \dot{\theta}_i z_i \\
v_i &= R_i^T (v_{i-1} + \omega_{i-1} \times p_i) + h \dot{\theta}_i z_i
\end{align}

\textbf{注意}：在这种情况下，旋转和移动是耦合的，都依赖于同一个关节变量$\dot{\theta}_i$。

\textbf{一般情况：独立的旋转和移动}

如果相对旋转和相对移动是独立的，那么：
\begin{itemize}
\item $\omega_{\text{rel}} = \dot{\theta}_i z_i$（绕轴$z_i$旋转）
\item $v_{\text{rel}} = \dot{d}_i z_i$（沿轴$z_i$移动，但与旋转独立）
\end{itemize}

因此：
\begin{align}
\omega_i &= R_i^T \omega_{i-1} + \dot{\theta}_i z_i \\
v_i &= R_i^T (v_{i-1} + \omega_{i-1} \times p_i) + \dot{d}_i z_i
\end{align}

\textbf{注意}：在这种情况下，旋转和移动是独立的，分别依赖于不同的关节变量$\dot{\theta}_i$和$\dot{d}_i$。

\textbf{最一般情况：任意相对运动}

如果相对旋转和相对移动沿着不同的轴，那么：
\begin{itemize}
\item $\omega_{\text{rel}} = \dot{\theta}_i z_i$（绕轴$z_i$旋转）
\item $v_{\text{rel}} = \dot{d}_j z_j$（沿轴$z_j$移动，$z_j \neq z_i$）
\end{itemize}

因此：
\begin{align}
\omega_i &= R_i^T \omega_{i-1} + \dot{\theta}_i z_i \\
v_i &= R_i^T (v_{i-1} + \omega_{i-1} \times p_i) + \dot{d}_j z_j
\end{align}

\textbf{注意}：在这种情况下，旋转和移动沿着不同的轴，完全独立。

\section{具体例子}

\subsection{例子1：螺旋关节（丝杠传动）}

\textbf{场景}：

考虑一个丝杠传动系统，其中：
\begin{itemize}
\item 丝杠的螺距为$h = 5$ mm/转
\item 丝杠以角速度$\dot{\theta} = 2\pi$ rad/s旋转
\item 螺母沿丝杠轴线移动
\end{itemize}

\textbf{分析}：

对于螺旋关节，旋转和移动是耦合的：
\begin{itemize}
\item $\omega_{\text{rel}} = \dot{\theta} z = 2\pi z$ rad/s（绕$z$轴旋转）
\item $v_{\text{rel}} = h \dot{\theta} z = 5 \times 10^{-3} \times 2\pi z = 0.01\pi z$ m/s（沿$z$轴移动）
\end{itemize}

\textbf{速度合成}：

假设坐标系$\{i-1\}$（丝杠固定端）的速度和角速度已知：
\begin{align}
\omega_{i-1} &= [0, 0, 0]^T \text{ rad/s} \\
v_{i-1} &= [0, 0, 0]^T \text{ m/s}
\end{align}

假设$R_i = I$（单位矩阵，即坐标系$\{i\}$与$\{i-1\}$对齐），$p_i = [0, 0, 0]^T$。

那么：
\begin{align}
\omega_i &= R_i^T \omega_{i-1} + \omega_{\text{rel}} = [0, 0, 0]^T + [0, 0, 2\pi]^T = [0, 0, 2\pi]^T \text{ rad/s} \\
v_i &= R_i^T (v_{i-1} + \omega_{i-1} \times p_i) + v_{\text{rel}} = [0, 0, 0]^T + [0, 0, 0.01\pi]^T = [0, 0, 0.01\pi]^T \text{ m/s}
\end{align}

\textbf{物理意义}：
\begin{itemize}
\item 螺母的角速度为$2\pi$ rad/s（1转/秒）
\item 螺母的线速度为$0.01\pi \approx 0.0314$ m/s（约3.14 cm/s）
\item 旋转和移动是耦合的：每转一圈，螺母移动$h = 5$ mm
\end{itemize}

\subsection{例子2：一般关节（旋转+独立移动）}

\textbf{场景}：

考虑一个复合关节，其中：
\begin{itemize}
\item 关节可以绕$z$轴旋转，角速度为$\dot{\theta} = 1$ rad/s
\item 关节可以沿$x$轴独立移动，速度为$\dot{d} = 0.1$ m/s
\item 旋转和移动是独立的（可以单独控制）
\end{itemize}

\textbf{分析}：

对于一般关节，旋转和移动是独立的：
\begin{itemize}
\item $\omega_{\text{rel}} = \dot{\theta} z = [0, 0, 1]^T$ rad/s（绕$z$轴旋转）
\item $v_{\text{rel}} = \dot{d} x = [0.1, 0, 0]^T$ m/s（沿$x$轴移动）
\end{itemize}

\textbf{速度合成}：

假设坐标系$\{i-1\}$的速度和角速度已知：
\begin{align}
\omega_{i-1} &= [0, 0.5, 0]^T \text{ rad/s} \\
v_{i-1} &= [0, 0, 0.2]^T \text{ m/s}
\end{align}

假设$R_i$是绕$z$轴旋转$90°$的旋转矩阵：
\begin{equation}
R_i = \begin{bmatrix}
0 & -1 & 0 \\
1 & 0 & 0 \\
0 & 0 & 1
\end{bmatrix}
\end{equation}

假设$p_i = [0, 0, 0.5]^T$ m。

\textbf{计算角速度}：
\begin{align}
R_i^T \omega_{i-1} &= \begin{bmatrix}
0 & 1 & 0 \\
-1 & 0 & 0 \\
0 & 0 & 1
\end{bmatrix} \begin{bmatrix}
0 \\ 0.5 \\ 0
\end{bmatrix} = \begin{bmatrix}
0.5 \\ 0 \\ 0
\end{bmatrix} \\
\omega_i &= R_i^T \omega_{i-1} + \omega_{\text{rel}} = [0.5, 0, 0]^T + [0, 0, 1]^T = [0.5, 0, 1]^T \text{ rad/s}
\end{align}

\textbf{计算线速度}：
\begin{align}
\omega_{i-1} \times p_i &= \begin{bmatrix}
0 \\ 0.5 \\ 0
\end{bmatrix} \times \begin{bmatrix}
0 \\ 0 \\ 0.5
\end{bmatrix} = \begin{bmatrix}
-0.25 \\ 0 \\ 0
\end{bmatrix} \\
v_{i-1} + \omega_{i-1} \times p_i &= [0, 0, 0.2]^T + [-0.25, 0, 0]^T = [-0.25, 0, 0.2]^T \\
R_i^T (v_{i-1} + \omega_{i-1} \times p_i) &= \begin{bmatrix}
0 & 1 & 0 \\
-1 & 0 & 0 \\
0 & 0 & 1
\end{bmatrix} \begin{bmatrix}
-0.25 \\ 0 \\ 0.2
\end{bmatrix} = \begin{bmatrix}
0 \\ 0.25 \\ 0.2
\end{bmatrix} \\
v_i &= R_i^T (v_{i-1} + \omega_{i-1} \times p_i) + v_{\text{rel}} = [0, 0.25, 0.2]^T + [0.1, 0, 0]^T = [0.1, 0.25, 0.2]^T \text{ m/s}
\end{align}

\textbf{物理意义}：
\begin{itemize}
\item 坐标系$\{i\}$的角速度为$[0.5, 0, 1]^T$ rad/s，包含：
\begin{itemize}
\item 从$\{i-1\}$传递过来的角速度：$[0.5, 0, 0]^T$ rad/s（在$\{i\}$中表示）
\item 相对旋转角速度：$[0, 0, 1]^T$ rad/s
\end{itemize}
\item 坐标系$\{i\}$的线速度为$[0.1, 0.25, 0.2]^T$ m/s，包含：
\begin{itemize}
\item 从$\{i-1\}$传递过来的速度：$[0, 0.25, 0.2]^T$ m/s（在$\{i\}$中表示）
\item 相对移动速度：$[0.1, 0, 0]^T$ m/s
\end{itemize}
\end{itemize}

\subsection{例子3：机器人手臂中的螺旋关节}

\textbf{场景}：

考虑一个6自由度机器人手臂，其中第3个关节是螺旋关节（用于精密定位）：
\begin{itemize}
\item 关节3的螺旋节距：$h_3 = 2$ mm/转
\item 关节3的角速度：$\dot{\theta}_3 = \pi$ rad/s（0.5转/秒）
\item 关节3的旋转轴：$z_3 = [0, 0, 1]^T$（在坐标系$\{3\}$中）
\end{itemize}

\textbf{前向递推计算}：

假设前两个关节的速度和角速度已知：
\begin{align}
\omega_2 &= [0, 0, 0.5]^T \text{ rad/s} \\
v_2 &= [0.1, 0, 0]^T \text{ m/s}
\end{align}

假设$R_3 = I$，$p_3 = [0, 0, 0.3]^T$ m。

\textbf{计算关节3的速度}：
\begin{align}
\omega_{\text{rel}} &= \dot{\theta}_3 z_3 = \pi \times [0, 0, 1]^T = [0, 0, \pi]^T \text{ rad/s} \\
v_{\text{rel}} &= h_3 \dot{\theta}_3 z_3 = 2 \times 10^{-3} \times \pi \times [0, 0, 1]^T = [0, 0, 0.002\pi]^T \text{ m/s}
\end{align}

\begin{align}
\omega_3 &= R_3^T \omega_2 + \omega_{\text{rel}} = [0, 0, 0.5]^T + [0, 0, \pi]^T = [0, 0, 0.5 + \pi]^T \text{ rad/s} \\
v_3 &= R_3^T (v_2 + \omega_2 \times p_3) + v_{\text{rel}}
\end{align}

计算$\omega_2 \times p_3$：
\begin{equation}
\omega_2 \times p_3 = \begin{bmatrix}
0 \\ 0 \\ 0.5
\end{bmatrix} \times \begin{bmatrix}
0 \\ 0 \\ 0.3
\end{bmatrix} = \begin{bmatrix}
0 \\ 0 \\ 0
\end{bmatrix}
\end{equation}

因此：
\begin{align}
v_3 &= R_3^T (v_2 + \omega_2 \times p_3) + v_{\text{rel}} \\
&= [0.1, 0, 0]^T + [0, 0, 0.002\pi]^T \\
&= [0.1, 0, 0.002\pi]^T \text{ m/s} \\
&\approx [0.1, 0, 0.0063]^T \text{ m/s}
\end{align}

\textbf{物理意义}：
\begin{itemize}
\item 关节3的角速度为$0.5 + \pi \approx 3.64$ rad/s
\item 关节3的线速度为$[0.1, 0, 0.0063]^T$ m/s
\item 由于螺旋关节的特性，每转一圈，关节3沿$z$轴移动$h_3 = 2$ mm
\item 在$\pi$ rad/s的角速度下，线速度为$0.002\pi \approx 0.0063$ m/s = 6.3 mm/s
\end{itemize}

\subsection{例子4：对比：旋转关节 vs 螺旋关节}

\textbf{场景}：

比较两个相似的关节，一个旋转关节，一个螺旋关节，都绕$z$轴旋转，角速度都是$\dot{\theta} = 2\pi$ rad/s。

\textbf{旋转关节}：
\begin{itemize}
\item $\omega_{\text{rel}} = 2\pi z$ rad/s
\item $v_{\text{rel}} = 0$
\item 结果：只有旋转，没有移动
\end{itemize}

\textbf{螺旋关节}（螺距$h = 5$ mm/转）：
\begin{itemize}
\item $\omega_{\text{rel}} = 2\pi z$ rad/s
\item $v_{\text{rel}} = h \times 2\pi z = 5 \times 10^{-3} \times 2\pi z = 0.01\pi z$ m/s
\item 结果：既有旋转，又有移动
\end{itemize}

\textbf{对比}：

假设$\omega_{i-1} = 0$，$v_{i-1} = 0$，$R_i = I$，$p_i = 0$。

\textbf{旋转关节}：
\begin{align}
\omega_i &= 0 + 2\pi z = [0, 0, 2\pi]^T \text{ rad/s} \\
v_i &= 0 + 0 = [0, 0, 0]^T \text{ m/s}
\end{align}

\textbf{螺旋关节}：
\begin{align}
\omega_i &= 0 + 2\pi z = [0, 0, 2\pi]^T \text{ rad/s} \\
v_i &= 0 + 0.01\pi z = [0, 0, 0.01\pi]^T \approx [0, 0, 0.0314]^T \text{ m/s}
\end{align}

\textbf{关键区别}：
\begin{itemize}
\item 旋转关节：只有角速度，没有线速度
\item 螺旋关节：既有角速度，又有线速度（沿旋转轴方向）
\item 螺旋关节的线速度与角速度成正比，比例系数是螺距$h$
\end{itemize}

\section{与经典力学的对应关系}

\subsection{经典力学中的速度合成定理}

在经典力学中，速度合成定理通常表述为：
\begin{equation}
\mathbf{v}_{\text{绝对}} = \mathbf{v}_{\text{牵连}} + \mathbf{v}_{\text{相对}}
\end{equation}

其中：
\begin{itemize}
\item $\mathbf{v}_{\text{绝对}}$：物体相对于固定参考系的速度
\item $\mathbf{v}_{\text{牵连}}$：由于参考坐标系运动产生的速度
\item $\mathbf{v}_{\text{相对}}$：物体相对于参考坐标系的速度
\end{itemize}

\subsection{对应关系}

在我们的公式中：
\begin{itemize}
\item $\mathbf{v}_{\text{绝对}} = v_i$（坐标系$\{i\}$原点的绝对速度）
\item $\mathbf{v}_{\text{牵连}} = R_i^T (v_{i-1} + \omega_{i-1} \times p_i)$（由于坐标系$\{i-1\}$运动产生的速度）
\item $\mathbf{v}_{\text{相对}} = v_{\text{rel}}$（坐标系$\{i\}$相对于坐标系$\{i-1\}$的速度）
\end{itemize}

\textbf{注意}：在机器人学中，我们还需要考虑坐标变换（$R_i^T$），因为不同坐标系中的表示需要转换。

\section{几何解释}

\subsection{角速度合成的几何解释}

\begin{itemize}
\item 角速度是矢量，可以按照矢量加法进行合成
\item 牵连角速度和相对角速度在同一个坐标系（$\{i\}$）中表示后，可以直接相加
\end{itemize}

\subsection{线速度合成的几何解释}

\begin{itemize}
\item 线速度合成需要考虑：
\begin{enumerate}
\item 基点（坐标系$\{i-1\}$原点）的速度
\item 由于旋转产生的速度（$\omega_{i-1} \times p_i$）
\item 相对速度
\end{enumerate}
\item 所有这些速度都需要在同一个坐标系（$\{i\}$）中表示后才能相加
\end{itemize}

\section{相对速度/角速度的表示约定}

\subsection{问题：$\omega_{\text{rel}}$和$v_{\text{rel}}$在哪个坐标系中表示？}

\textbf{问题}：如果$\omega_{\text{rel}}$是坐标系$\{i\}$相对于坐标系$\{i-1\}$的角速度，那么它应该在哪个坐标系中表示？

\textbf{答案}：这取决于定义约定。在机器人学中，通常的约定是：

\subsection{约定1：在坐标系$\{i\}$中表示（常用）}

\textbf{原因}：
\begin{itemize}
\item 关节轴$z_i$通常在坐标系$\{i\}$中定义
\item 因此$\omega_{\text{rel}} = \dot{\theta}_i z_i$自然在$\{i\}$中表示
\item 同样，$v_{\text{rel}} = \dot{d}_i z_i$也在$\{i\}$中表示
\end{itemize}

\textbf{优点}：
\begin{itemize}
\item 与关节轴的定义一致
\item 在递归算法中，所有量都在当前坐标系$\{i\}$中表示，便于计算
\end{itemize}

\textbf{公式}：
\begin{align}
\omega_i &= R_i^T \omega_{i-1} + \omega_{\text{rel}} \tag{3} \\
v_i &= R_i^T (v_{i-1} + \omega_{i-1} \times p_i) + v_{\text{rel}} \tag{4}
\end{align}

其中$\omega_{\text{rel}}$和$v_{\text{rel}}$都在坐标系$\{i\}$中表示。

\subsection{约定2：在坐标系$\{i-1\}$中表示（理论可能）}

\textbf{理论上}，相对角速度也可以在$\{i-1\}$中表示：
\begin{equation}
\omega_{\text{rel},\{i-1\}} = R_i \omega_{\text{rel},\{i\}} = R_i (\dot{\theta}_i z_i)
\end{equation}

\textbf{但这样会导致公式更复杂}：
\begin{equation}
\omega_i = R_i^T (\omega_{i-1} + \omega_{\text{rel},\{i-1\}})
\end{equation}

\textbf{缺点}：
\begin{itemize}
\item 需要额外的坐标变换
\item 与关节轴的定义不一致
\item 在递归算法中不常用
\end{itemize}

\subsection{标准约定（机器人学）}

在机器人学中，\textbf{标准约定}是：
\begin{itemize}
\item $\omega_{\text{rel}}$和$v_{\text{rel}}$都在坐标系$\{i\}$中表示
\item 这样，公式(3)和(4)中的各项都在同一个坐标系$\{i\}$中，可以直接相加
\item 这与关节轴$z_i$在$\{i\}$中定义一致
\end{itemize}

\textbf{物理意义}：
\begin{itemize}
\item $\omega_{\text{rel}}$是$\{i\}$相对于$\{i-1\}$的角速度
\item 虽然它是"相对于$\{i-1\}$"的，但它的\textbf{表示}在$\{i\}$中
\item 这类似于：虽然我们说"相对于地面"，但速度可以在任何坐标系中表示
\end{itemize}

\subsection{类比说明}

\textbf{类比}：考虑一个在地面上行驶的汽车
\begin{itemize}
\item 汽车的速度是"相对于地面"的
\item 但这个速度可以在以下坐标系中表示：
\begin{itemize}
\item 地面坐标系（固定坐标系）
\item 汽车坐标系（移动坐标系）
\item 其他任意坐标系
\end{itemize}
\item 虽然速度是"相对于地面"的，但它的表示可以在任何坐标系中
\end{itemize}

\textbf{同样}：
\begin{itemize}
\item $\omega_{\text{rel}}$是$\{i\}$相对于$\{i-1\}$的角速度
\item 但这个角速度可以在$\{i\}$中表示（标准约定）
\item 也可以在$\{i-1\}$中表示（需要坐标变换）
\end{itemize}

\section{总结}

\subsection{关键要点}

\begin{enumerate}
\item \textbf{官方名称}：
\begin{itemize}
\item 角速度合成定理（Angular Velocity Composition Theorem）
\item 线速度合成定理（Linear Velocity Composition Theorem）
\item 或统称为刚体速度合成定理（Rigid Body Velocity Composition Theorem）
\end{itemize}

\item \textbf{物理意义}：
\begin{itemize}
\item 绝对速度/角速度 = 牵连速度/角速度 + 相对速度/角速度
\item 牵连部分包括平移和旋转两部分
\end{itemize}

\item \textbf{数学表示（一般形式）}：
\begin{align}
\omega_i &= R_i^T \omega_{i-1} + \omega_{\text{rel}} \tag{3} \\
v_i &= R_i^T (v_{i-1} + \omega_{i-1} \times p_i) + v_{\text{rel}} \tag{4}
\end{align}

\item \textbf{关键点}：
\begin{itemize}
\item 所有速度/角速度都需要在同一个坐标系中表示后才能相加
\item 需要进行坐标变换（$R_i^T$）
\item 线速度合成需要考虑旋转产生的速度（$\omega_{i-1} \times p_i$）
\item 公式(3)和(4)适用于所有情况，包括：
\begin{itemize}
\item 只有旋转（$\omega_{\text{rel}} \neq 0$，$v_{\text{rel}} = 0$）
\item 只有移动（$\omega_{\text{rel}} = 0$，$v_{\text{rel}} \neq 0$）
\item 既有旋转又有移动（$\omega_{\text{rel}} \neq 0$，$v_{\text{rel}} \neq 0$）
\end{itemize}
\end{itemize}
\end{enumerate}

\subsection{三种情况的对比}

\begin{table}[h]
\centering
\begin{tabular}{|l|c|c|c|}
\hline
\textbf{关节类型} & $\omega_{\text{rel}}$ & $v_{\text{rel}}$ & \textbf{典型应用} \\
\hline
旋转关节 & $\dot{\theta}_i z_i$ & $0$ & 大多数机器人关节 \\
\hline
移动关节 & $0$ & $\dot{d}_i z_i$ & 线性驱动器、伸缩臂 \\
\hline
螺旋关节 & $\dot{\theta}_i z_i$ & $h \dot{\theta}_i z_i$ & 螺旋传动、丝杠 \\
\hline
一般关节 & $\dot{\theta}_i z_i$ & $\dot{d}_j z_j$ & 复合关节、特殊机构 \\
\hline
\end{tabular}
\caption{不同关节类型的相对速度/角速度}
\end{table}

\textbf{注意}：
\begin{itemize}
\item 对于旋转关节和移动关节，公式(3)和(4)中的$\omega_{\text{rel}}$或$v_{\text{rel}}$为零
\item 对于螺旋关节，$\omega_{\text{rel}}$和$v_{\text{rel}}$是耦合的（都依赖于$\dot{\theta}_i$）
\item 对于一般关节，$\omega_{\text{rel}}$和$v_{\text{rel}}$可以完全独立
\end{itemize}

\subsection{应用}

这些公式在机器人学中用于：
\begin{itemize}
\item 递归计算每个连杆的速度和角速度
\item 前向运动学计算
\item 雅可比矩阵的计算
\item 动力学分析
\end{itemize}

\end{document}

